\cotefile{149}

\scanpage{103}
\mspage{51}
On a évidemment $\mathrm{a)}\Rightarrow\mathrm{b)}\Rightarrow\mathrm{c)}\unsure{=|\Rightarrow}\mathrm{d)}$, et $d\Rightarrow a$ devrait être plus ou moins connu, ou du moins se prouver à l'aide de choses bien connues. Ces conditions impliquent e, mais je n'ai pas d'idée pour le moment comment voir que e) implique bonne réduction --- faudrait une compréhension (fut-elle heuristique) plus approfondie que celle que j'ai pour le moment.
\margin{Ça vaudrait la peine de regarder une version de ceci où $X_\eta=\hat U_\eta$ a bonne réduction --- mais la multisection $\overline{X_\eta-U_\eta}$ de $X$ sur $S$ n'est pas étale. Le cas du genre $0$ devrait être particulièrement aisé à regarder\ldots}

J'ai donc là un principe général (que je pourrais à l'occasion formuler avec la généralité qui lui convient) pour construire des actions extérieures \inserted{Individuelles} de groupes $\mathcal E_k$ ($k$ \unsure{ext.|\deleted{ext.}} de type fini de $\mathbf Q$) sur des $\pi$ à \unsure{bouts|lacets}, qui \inserted{(sans doute)} ne soient géométriques, par \deleted{des} la considération de courbes de type $g,\nu$ ``se réduisant mal''. Mais je n'arrive
\margin{102}

\scanpage{105}
\mspage{52}
4. \uline{Sections} \inserted{d'extensions,} \uline{et anneaux de valuation généraux.}

D'abord une révision de notations. Si $X$ est un schéma connexe, je \deleted{\unsure{dé}}note \deleted{$\mathcal E_X$}
\[
(1)\qquad \mathcal E_X=\pi_1(X)
\]
son groupe fond. profini en tant que groupe extérieur, et si $\tilde X$ est un rev. universel profini de $\mathcal E_X$, par
\[
(2)\qquad \mathcal E_X^{(\tilde X)}=\pi_1(X;\tilde X)=\operatorname{Aut}_X(\tilde X)
\]
son groupe fondamental précisé --- qui est un groupe profini. Si $X=\operatorname{Spec}(A)$, où $A$ est un anneau (le plus souvent un corps) je note $\mathcal E_A$, et $\mathcal E_A^{\tilde A}$ \inserted{$=\mathcal E_X^{\unsure{X|\tilde X}}$} si $\unsure{\tilde A|A}$ est une $A$-algèbre telle que $\operatorname{Spec}(\tilde A)=\tilde X$ soit un rev. universel de $X$ (ce qui le détermine à isom non unique près). Bien entendu, si $\unsure{\xi|\zeta}$ est un ``point géométrique'' de $X$, on note
\[
(3)\qquad \mathcal E_X^{\unsure{\xi|\zeta}}=\pi_1(X,\unsure{\xi|\zeta})=\mathcal E_X^{\tilde X(\unsure{\xi|\zeta})},
\]
où $\tilde X(\unsure{\xi|\zeta})$ est le rev universel de $X$ pointé en $\unsure{\xi|\zeta}$. Le choix
\margin{104}

\scanpage{106}
de $\unsure{\xi|\zeta}$ correspond au choix d'un $x\in X$ et d'une ext. sép. close \deleted{$\bar K$} \inserted{$\Omega$} de $k(x)$ (pas néc. \deleted{\unsure{une}} clôture alg. de $k(x)$) et on note alors aussi $\mathcal E_X^\Omega$ au lieu de $\mathcal E_X^{\unsure{\xi|\zeta}}$ ($\Omega$ sous-entendu en tant qu'extension de $k(x)$, donc avec sa structure de $k(x)$ algèbre). On aura \inserted{un isom can}
\[
(4)\qquad \mathcal E_x^\Omega\simeq\mathcal E_X^{\overline{k(x)}},
\]
où $\overline{k(x)}$ est la clôture \unsure{d|s} alg. de $k(x)$ dans $\Omega$. Bien sûr, si $X=\operatorname{Spec}(A)$, on notera aussi $\mathcal E_A^\Omega$ --- notation surtout utilisée pour $\mathcal E_K^{\bar K}$, $K$ un corps, $\bar K$ une clôture alg. ou sép. de $K$.

Si $X$ est un $Y$-schéma, \inserted{$Y$ connexe,} on \unsure{notera|\deleted{notera}}
\[
(5)\qquad \mathcal E(f):\mathcal E_X\longrightarrow\mathcal E_Y\qquad\text{où }f:X\longrightarrow Y
\]
($\mathcal E_X$ est un foncteur en $X$
\[
(\text{Sch conn})\longrightarrow\text{grpes ext}
\]
\margin{105}

\scanpage{107}
\mspage{53}

qui se précise par l'hom effectif de groupes profinis
\[
(6)\qquad \mathcal E_X^{\widetilde X}\longrightarrow\mathcal E_Y^{\widetilde Y}
\]
où $\widetilde Y$ est le rev. universel de $Y$ défini par $\widetilde X\longrightarrow Y$ ($\widetilde X$ jouant le rôle du point --- on devrait pouvoir écrire en fait $\mathcal E_Y^{\widetilde X}$, plus gen. $\mathcal E_Y^Z$ chaque fois qu'on a un $Y$-schéma $Z$ $1$-connexe, jouant le rôle de ``foncteur fibre'' pour le topos $B_{\pi(X)}$ des rev étales de $Y$\ldots)

\deleted{\unsure{que} $\mathcal E_X^{\widetilde{\unsure{X|\alpha}}}$}

On peut dire que
\[
(7)\qquad (X,\widetilde X)\longmapsto\mathcal E_X^{\widetilde X}
\]
est un foncteur, de la catégorie des schémas $0$-connexes $X$ \uline{munis} d'un rev. universel (ou seulement d'un $Z$ $1$-connexe s'envoyant dans $X$) vers celle des groupes profinis. Ceci s'applique en particulier en regardant la \inserted{sous-}cat. des $(X,\unsure{\xi|\zeta})$ munis d'un $p^t$ géom ---

\margin{106}

\scanpage{108}

on a donc
\[
(8)\qquad \mathcal E_X^{\unsure{\xi|\zeta}}\longrightarrow\mathcal E_{\deleted{\unsure{X|Y}}Y}^{\eta}
\]
si on a un hom de schémas ``géom\textsuperscript{t}. ponctués'' $(X,\unsure{\xi|\zeta})\longrightarrow(Y,\eta)$. On notera que tte ponctuation géom. de $X$ en un $x\in X$ --- i.e. une ext. sep close $\Omega$ de $k(x)$ --- définit \inserted{can\textsuperscript{t}} une ponctuation \inserted{géom} aussi de $Y$ en $y=f(x)$, \deleted{par} par composition --- $\Omega$ ``est'' une extension de $k(y)$ puisque $k(\underset{\inserted{x}}{x})\hookleftarrow k(y)$ --- et l'hom (8) s'identifie aussi à
\[
(9)\qquad \mathcal E_X^{\overline{k(x)}}\longrightarrow\mathcal E_Y^{\overline{k(y)}},
\]
où $\overline{k(x)}$, $\overline{k(y)}$ sont les clôtures séparables \uline{dans $\Omega$}.

On posera
\[
(10)\qquad \mathcal E_{X/Y}^{\widetilde X}=\operatorname{Ker}\bigl(\mathcal E_X^{\widetilde X}\longrightarrow\mathcal E_Y^{\widetilde X}=\mathcal E_Y^{\widetilde Y}\bigr)
\]

\margin{107}

\scanpage{109}
\mspage{54}

(C'est un foncteur p.r. au triple $\widetilde X\mathrel{\text{---}}X\mathrel{\text{---}}Y$ avec $X,Y$ $0$-connexes, $\widetilde X$ un rev. universel; plus généralt, si $\deleted{Z}\inserted{T}\longrightarrow X$ avec $\deleted{Z}\inserted{T}$ $1$-connexe, on pose
\[
(11)\qquad \unsure{\mathcal E}_{X/Y}^{\deleted{Z}\inserted{T}}=\operatorname{Ker}\bigl(\mathcal E_X^{\deleted{Z}\inserted{T}}\longrightarrow\mathcal E_Y^{\deleted{Z}\inserted{T}}\bigr)
\]
on a un foncteur en \inserted{le triple} $(\deleted{Z}\inserted{T}\longrightarrow\unsure{X_0|X}\longrightarrow Y)$ variable. En particulier: $\mathcal E_{X/Y}^{\unsure{\xi|\zeta}}$, $\unsure{\xi|\zeta}$ un $p^t$ géom. de $X$, $\mathcal E_{X/Y}^{\Omega}$, $\mathcal E_{X/A}^{\widetilde X}$ (si $Y=\operatorname{Spec} A$), $\mathcal E_{B/A}^{\widetilde B}$ $\ldots$

\margin{\uline{NB} $\mathcal E_{X/Y}^T$, pour $T$ variable, $X/Y$ fixé, est défini à isom non unique ! près, et se note $\mathcal E_{X/Y}$ simplem\textsuperscript{t}.}

Un cas intéressant est celui où on dispose d'une ``suite exacte d'homotopie universelle'' \deleted{de} (en dim $1$) pour $X\longrightarrow Y$, alors la donnée (pour $X=Y$ donné) de $T\longrightarrow X$, (avec $T$ $1$-connexe) peut s'interpréter par la donnée du composé

\margin{Car $\mathcal E_X^{\deleted{Z}\inserted{T}}\longrightarrow\mathcal E_Y^{\deleted{Z}\inserted{T}}$ est épi}
\margin{108}

\scanpage{110}

\[
T\longrightarrow Y
\]
\uline{et} d'un relèvement en $T\longrightarrow X$, ou ce qui revient au $=$, d'une section de $X_T=X\times_YT$ sur $T$. Ceci \unsure{posé,|pose;}

\noindent
\begin{minipage}[t]{.45\linewidth}
\vspace{\baselineskip}
\unsure{\diagram{\begin{tikzcd}
X\arrow[d] & X_T\arrow[l]\arrow[d] & \\
Y & T\arrow[l] & T\arrow[ul,leftrightarrow]\arrow[l,phantom,"=" description]
\end{tikzcd}}}
\end{minipage}\hfill
\begin{minipage}[t]{.51\linewidth}
\vspace{0pt}
on aura un isom can (sous l'hyp. de ``suite exacte d'homotopie'' faite)
\end{minipage}

\[
(12)\qquad \mathcal E_{X/Y}^T\simeq\pi_1(X_T;T)\simeq\mathcal E_{X_T}^T
\]
et on aura
\[
(13)\qquad 1\longrightarrow\mathcal E_{X/Y}^T\longrightarrow\mathcal E_X^T\longrightarrow\mathcal E_Y^T\longrightarrow1
\]
(Cette hypothèse sera tjrs satisfaite si $Y=\operatorname{Spec}K$, $K$ un corps, si $X$ est géom. $0$-connexe sur $K$)

Plus généralement, si on a une suite exacte d'homotopie universelle, mais pas néc. avec $\mathcal E_X^T\longrightarrow\mathcal E_Y^T$ surj,

\margin{109}

\scanpage{111}
\mspage{55}

\margin{\unsure{Sous|dans} l'hyp. ``suite exacte d'homo-\newline topie'', mais avec fibres pas néc. \unsure{conn|conn.}}

on aura une factorisation can. de $X\longrightarrow Y$ en
\[
(14)\qquad X\longrightarrow Y'\longrightarrow Y
\]
avec $Y'$ étale fini ou pro-étale fini sur $Y$ et $\mathcal E_X'\longrightarrow\mathcal E_{Y'}$ est maintenant épi, \deleted{\unsure{donc}} cette suite exacte univ. d'homotopie bien sûr. On aura des iso can
\[
(15)\qquad \mathcal E_{X/Y}^T\xrightarrow{\sim}\mathcal E_{X/Y'}^T
\]
qui peut donc se calculer comme
\[
(16)\qquad \mathcal E_{X\times_{Y'}T}^T\simeq\mathcal E_{X/Y}^T
\]
qu'on peut noter $\mathcal E_{X_T}^T$, mais en faisant attention que maintenant $X_T$ désigne non plus $X\times_YT$ (qui va être disconnexe \inserted{si $Y'\longrightarrow Y$ pas iso}) mais $X\times_YT$.

Bien sûr, à iso. can. pres on peut \deleted{\unsure{en}} remplacer, pour le calcul de $\mathcal E_{X_{\deleted{\unsure{X'|X}}\inserted{T}}}^T$, $T$ par $\widetilde Y$ pour ce qui est du $X_T$, et par $\widetilde X$ pour ce qui est de l'ancrage dans $X_{\widetilde Y}$,

\margin{110}

\scanpage{112}

où $\widetilde X,\widetilde Y$ sont les rev. universels de $X,Y$ déduits de \inserted{$\deleted{\unsure{X}}T\unsure{\mathrel{\text{---}}|\longrightarrow}X$ et de $T\unsure{\mathrel{\text{---}}|\longrightarrow}Y$,} de telle façon que l'on obtient
\[
(17)\qquad \mathcal E_{X/Y}^T\simeq\mathcal E_{X/Y}^{\widetilde X}\simeq\operatorname{Ker}\bigl(\mathcal E_X^{\widetilde X}\unsure{\mathrel{\text{---}}|\longrightarrow}\mathcal E_Y^{\widetilde Y}\bigr)
\]
(sans hyp de suite exacte)

et
\[
\begin{aligned}
(18)\qquad \mathcal E_{X/Y}^T\simeq{}&\deleted{\operatorname{Ker}\bigl(\unsure{\mathcal E_{X_{\widetilde Y}}^{\widetilde X}}\unsure{\mathrel{\text{---}}|\longrightarrow}\unsure{\mathcal E_{\widetilde Y}^{\widetilde Y}}\bigr)}\\
&\inserted{\mathcal E_{X_{\widetilde Y}}^Z\simeq\mathcal E_{X_{\widetilde Y}}^{\widetilde X}}
\end{aligned}
\]
dans le cas général.

\diagram{\begin{tikzcd}
X\arrow[d] & X_{\widetilde Y}\arrow[l]\arrow[d] & \widetilde X\arrow[l] & Z\arrow[l]\\
Y'\arrow[d] & \widetilde Y\arrow[l] & & \\
Y & & &
\end{tikzcd}}

Une autre interprétation de $\mathcal E_{X/Y}^T$ s'obtient en prenant $T$ réduit à un $p^t$ géométrique \inserted{$\unsure{\xi|\zeta}_X$}, qui définit donc \unsure{un|une} $p^t$ géom. $\unsure{\xi|\zeta}_Y$ de $Y$, une fibre géom. $X_{\unsure{\xi|\zeta}_Y}$ de $X$ sur $Y$ en $\unsure{\xi|\zeta}_Y$, $\unsure{\xi|\zeta}_x$ s'interprétant comme un $p^t$ géom. de cette fibre, et on aura
\[
(19)\qquad \mathcal E_{X/Y}^{\deleted{\unsure{T}}\inserted{\unsure{\xi|\zeta}_X}}\simeq\mathcal E_{X_{(\unsure{\xi|\zeta}_Y)}}^{\unsure{\xi|\zeta}_X}\simeq\pi_1(X_{(\unsure{\xi|\zeta}_Y)},\unsure{\xi|\zeta}_X)
\]
--- i.e. $\mathcal E_{X/Y}$ s'interprète comme groupe fondamental des ``fibres géom'' de $X/Y$.

\margin{111}

\scanpage{113}
\mspage{56}

Le lien avec la \deleted{\unsure{définition}}\inserted{description} (11) est évidemment que, si $\unsure{\xi|\zeta}_X$ provient d'un \unsure{pt|$p^t$} géom $\unsure{\xi|\zeta}_T$ de $T$, alors on aura des iso can

\diagram{\begin{tikzcd}[column sep=large,row sep=large]
(20) & \mathcal E_{X/Y}^{\unsure{\xi|\zeta}_T}\arrow[r,phantom,"\simeq" description]\arrow[d,phantom,"\rotatebox{90}{$\simeq$}" description] & \mathcal E_{X/Y}^T\\
& \pi_1(X_{\unsure{\xi|\zeta}_Y},\unsure{\xi|\zeta}_T) &
\end{tikzcd}}

\deleted{\unsure{donc|les}} --- les $\pi_1(X_{\unsure{\xi|\zeta}_Y},\unsure{\xi|\zeta}_T)$ des fibres géom, relatifs aux divers pts géom de $\mathcal E_{X/Y}^T$, sont canoniquement iso \unsure{entre eux.|entre: eux.}

( Il faut bien faire attention que $\mathcal E_{X/Y}$ \inserted{($Z$ non précisé)} est une classe d'iso. de groupes profinis sans plus, \uline{pas} un groupe extérieur. L'arbitraire dans le choix d'un iso entre $\mathcal E_{X/Y}^T$ et $\mathcal E_{X/Y}^{T'}$ provient du choix

\margin{112}

\scanpage{114}

d'un \unsure{él.|el.} de $\mathcal E_{X\deleted{\unsure{/Y}}}^T$, opérant sur le sous-groupe $\mathcal E_{X/Y}^T$ par autom. intérieurs.

Les notations \inserted{$\mathcal E_X$,} $\mathcal E_{X/Y}$ evoquent intuitivement un $p^t$ de vue qui regarde $\mathcal E_{X\deleted{\unsure{/Y}}}^T$ comme une \uline{extension} de $\mathcal E_Y^T$ (ou d'un ss groupe convenable seulement) par $\mathcal E_{X/Y}^T$, qui sera souvent, pour les besoins de la cause, noté $\pi$ --- il joue le rôle d'un groupe fondamental (géométrique), tandis que $\mathcal E_Y^T$ joue le rôle d'un groupe de Galois --- ou d'un groupe profini mixte.

\margin{et j'hésite encore entre les \unsure{notations|notation} $\mathcal E_{X/Y}^T$ et $\pi_{X/Y}^T$; que je vais peut-être employer concur-\newline remment.}

Le cas le plus important pour nous sera celui où $Y=\operatorname{Spec}K$, $K$ un corps. \deleted{\unsure{Aussi} \unsure{formellement|souvent}} Dans ce cas, $\mathcal E_Y\simeq\mathcal E_Y^{\widetilde Y}$

\margin{113}

\scanpage{115}
\mspage{57}

s'identifie bel et bien au groupe de Galois $\operatorname{Gal}(\overline K/K)$ de \deleted{$\unsure{\mathcal E_Y}$} $\overline K/K$, $\overline K$ étant la cloture sép. de $K$ telle que $\operatorname{Spec}\overline K\simeq\widetilde Y$. Souvent, on notera $\Gamma$, \inserted{ou $\Gamma_K$, $\Gamma_K^{\overline K}$,} au lieu de $\mathcal E_Y$ --- surtout si $K$ est alg. sur le corps premier, et qu'il n'a donc plus de ``partie géométrique'' à distinguer d'une ``partie arithmétique''$\ldots$

Soit \inserted{un} $K$ \deleted{\unsure{tel}} corps (qui pourrait être alg. clos), $L$ une ext. de type fini de $K$, $X$ un ``modele'' propre régulier de $L$. Alors $\mathcal E_X^{\overline L}$ s'identifie à un quotient de $\mathcal E_L^{\overline L}$, qui \uline{ne \deleted{\unsure{p}} dépend pas des \unsure{modèles|modèle} $X$ \unsure{choisi|choisis}}, comme il est bien connu

\margin{114}

\scanpage{116}

c'est la partie ``\unsure{universellement|universelles} géom'' de $\mathcal E_{\deleted{\unsure{K}}\inserted{L}}^{\overline L}$, \deleted{\unsure{noté} $\mathcal E_{L/K}^{\unsure{g|q}\overline L}$} \deleted{qui \unsure{couvre} les \unsure{hom|bons}} \inserted{(relativement à $K$)} qui classifie les \unsure{schémas|sch.} (finis) etales sur $L$ qui sont ``non ramifiés'' sur \uline{\unsure{tt|tout}} modèle \deleted{\unsure{étal}} régulier (propre ou non) de $L/K$.

\margin{noté $\mathcal E_L^{\unsure{g|q}\overline L}$ ou $\mathcal E_L^{\unsure{g(K)},\overline L}$}

Si $U$ est un modèle quelconque; il se plonge dans un $X$, et on a des hom surjectifs
\unsure{\deleted{$q$}|\deleted{$\mathcal E$}|$\mathcal E$}
\[
(21)\qquad \mathcal E_L^{\overline L}\longrightarrow\mathcal E_U^{\overline L}\longrightarrow\mathcal E_X^{\overline L}=\mathcal E_{\deleted{\unsure{X}}\inserted{L}}^{\unsure{g|q}\overline L}
\]

On aura \deleted{\unsure{alors}}
\[
(22)\qquad U=X\setminus Z
\]
$Z$ partie ferm. de $X$ --- soient $D_i$ ses composantes de codim $1$, $x_i$ les pt gén. de $D_i$. \unsure{\deleted{A}\inserted{O}|O}n a des groupes de décomposition et d'inertie $N_{x_i'}$ et $I_{x_i'}$ \inserted{\unsure{de|dans} $\mathcal E_U^{\overline L}$} \unsure{\deleted{d}|d} relatifs au choix d'un $x_i'$ au dessus de $x_i$ dans le \unsure{normalisé} de $X$ dans $\overline L$), et on trouve

\margin{115}

\scanpage{117}
\mspage{58}
\[
(23)\qquad
\mathcal E_X^{\deleted{\unsure{K}}\inserted{\overline L}}
=\mathcal E_L^{\unsure{g|G}\overline L}
\simeq\mathcal E_U^{\overline L}/\begin{gathered}\text{ss-groupe fermé}\\\text{inv. engendré}\\\text{par les }\unsure{L_{x_i'}|I_{x_i'}}\\(i\in I)\end{gathered}
\]
Notons que $\mathcal E_L^{\overline L}$ est la limite projective des $\mathcal E_U^{\overline L}$, pour des modèles réguliers (affines si on y tient) variables
\[
(24)\qquad\mathcal E_L^{\overline L}\simeq\lim_{\longleftarrow\,U}\mathcal E_U^{\overline L}
\]
et de $=$, bien sûr
\[
(25)\qquad\pi_{L/K}^{\overline L}\simeq\lim_{\longleftarrow\,U}\pi_{U/K}^{\overline L}.
\]
Comment récupérer les $\mathcal E_U^{\overline L}$ \unsure{à} $\mathcal E_L^{\overline L}$ ? Pour \uline{tt} sous-anneau \inserted{d-vr $V$} de \unsure{ess.|en.} de type \unsure{fini} sur $K$ (correspond à un ``germe de modèle'' de $L/K$, au voisinage d'un pt de codim $1$ i.e. du pt gén. d'un diviseur), considé-\newline
rons le point fermé $x$ de $S_x=\operatorname{Spec}(V)$,
\margin{Dans les considérations qui suivent, $K$ (et $X,U$) ne jouent provisoirement aucun rôle --- \unsure{(|[}Mais on pourrait ``plus gén.'' regarder groupes de décomp. \deleted{\unsure{de dimension relative}} d'un des ``$U$'' \uline{mais seulement} \uline{\unsure{épis-comp|épais-comp}} (la ``généralisation'' étant triviale, \unsure{par} $\mathcal E_L\longrightarrow\mathcal E_U$)}
\margin{NB on n'utilise \unsure{les|en} termes ``absolus'' que pour des schémas \uline{épais} (quasi-proj. ?) \uuline{réguliers}}
\margin{116}

\scanpage{118}
et un $x'$ au dessus de $x$ dans le normalisé de $S_x$ dans $\overline L$, d'où un groupe de décomposition $N_{x'}$, et un groupe d'inertie $L_{x'}$, avec
\[
(26)\qquad N_{x'}/I_{x'}\simeq k(x')
\]
(ext. de type fini de $K$, de degré de tr. $n-1$, si $n=\deg\operatorname{tr}L/K$), et
\[
(27)\qquad I_{x'}\simeq\unsure{T_0|T_\circ|T_\infty}(\overline L)\simeq T_\infty(\overline K)\;-
\]
(par Kummer) \deleted{Notons \unsure{Ainsi|bien}, \unsure{on a}} \inserted{Ces constructions} ne dependent du choix d'aucun \deleted{\unsure{de|des}} modele ($U$, ou $X\ldots$). --- Notons que
\[
(28)\qquad I_{x'}\subset\pi_{L/K}^{\overline L}
\]
En fait, on peut dire, posant
\[
(29)\qquad\left\{\begin{aligned}
\widetilde{\unsure{\deleted{U}\inserted{S}|S}}_x&=\operatorname{Spec}\widetilde V,\quad\widetilde V\text{ hensélisé (pas strict) de }V\\
\widetilde U_x&=\widetilde S_x\setminus\{x\}\simeq\operatorname{Spec}\widetilde{\deleted{K}\inserted{L}}_x\quad(\text{où }\widetilde{\deleted{K}\inserted{L}}_x\text{ est le corps des fractions de }\widetilde V)
\end{aligned}\right.
\]
que l'on a un hom \inserted{ext} $\mathcal E_{\widetilde U_x}\longrightarrow\mathcal E_L$, \deleted{(\unsure{alors})} dont le choix ``effectif'' dépend de
\margin{$I_{x'}$ joue \unsure{le} role d'un ``groupe bout''}
\margin{si Car. $k(x)$ est nulle --- sinon quotient \unsure{du seul habituel}\ldots}
\margin{117}

\scanpage{119}
\mspage{59}
celui d'un rev. universel \deleted{\unsure{(précis.}} \inserted{ou encore} d'un pt géom $\unsure{\xi|\zeta}$ de $\widetilde K_x$ --- i.e. \unsure{d'une|de'une} clôture algébrique de $\widetilde K_x$ --- \uline{et} d'une \deleted{\unsure{isom}} ``chemin de domaines'' entre $\unsure{\xi|\zeta}_{\widetilde U_x}\unsure{\mathrel{\text{---}}|\longrightarrow}\unsure{\widetilde K_x|\widetilde U_x}\longrightarrow\operatorname{Spec}(L)$, et $\unsure{\xi|\zeta}_L$ défini par $\widetilde L$ --- ou ce qui revient au $=$, \inserted{à isom près des \unsure{revts|relvts}} d'un $\deleted{\unsure{K}}\inserted{L}$-hom de \deleted{$\unsure{K_x|K}$} \deleted{$\unsure{K_x|K}$} $\widetilde L_x$ dans $\overline L$ --- et ceci à son tour équivaut à la donnée d'un $x'$, \deleted{\uline{et} \unsure{d'un is}} comme dessus. Ceci posé, on a des isom can.
\[
(30)\quad\left\{\begin{aligned}
N_{x'}&\simeq\mathcal E_{\widetilde U_x}^{\overline L}\\
I_{x'}&\simeq\text{Noyau de l'hom \unsure{canoniq|canonique} }\inserted{\text{surjectif}}\\
&\mathcal E_{\widetilde U_x}^{\overline L}\longrightarrow\mathcal E_{\overline{k(x)}}
\end{aligned}\right.
\]
(où la clôture \inserted{alg} $\overline{k(x)}$ de $k(x)$ est induite par la clôture alg $\overline L$ de $\widetilde K_x$)
\margin{Cf théorie galoisienne des traits henséliens Ici $I_{x'}$ joue le rôle d'un $\pi_{\widetilde U_x/k(x)}^{\overline L}$ ; en fait il est $=\pi_{\widetilde U_x/\widetilde S_x}^{\overline L}$ !}
(Alors (27) est déduit par théorie de Kummer)
\deleted{$I_{x'}\unsure{\simeq|=}$}
(\deleted{\unsure{27|28}})

\margin{? $\Vert$}
Je présume que (dans la situation actuelle --- ou même, avec un corps de base $K$) on a $N_{x'}=\operatorname{Norm}_{\mathcal E_L}(I_{x'})$, \inserted{(a priori, on a seulement $\subset$.)} ce qui est lié à l'idée qu'on doit avoir $I_{x'}=\operatorname{Norm}_{\pi_{L/K}}(I_{x'})$
\margin{118}

\scanpage{120}
--- à tirer au clair à l'occasion !

Revenant au modèle $U$, fournissant un quotient $\mathcal E_U^{\overline L}$ de $\mathcal E_L^{\overline L}$, d'où $N_{x'}\longrightarrow\mathcal E_U^{\overline L}$, induisant $I_{x'}\longrightarrow\mathcal E_U^{\overline L}$,
\deleted{Il est clair que \unsure{si}}
qui n'est autre que l'hom \deleted{\unsure{can}} $\mathcal E_{\widetilde U_x}\unsure{\mathrel{\text{---}}|\longrightarrow}\mathcal E_U$ associé à $\widetilde U_x\longrightarrow U$. \uline{Si} cet hom factorise par $\widetilde S_x$, il est immédiat que l'hom induit $I_{x'}\longrightarrow\mathcal E_U^{\overline L}$ est nul, puisque \inserted{$I_{x'}$} \deleted{\unsure{c'}}est le noyau de $\mathcal E_{\widetilde U_x}\Longrightarrow\mathcal E_{\widetilde S_x}$, et $\mathcal E_{\widetilde U_x}\unsure{\mathrel{\text{---}}|\longrightarrow}\mathcal E_U$ se factorise par cet hom.
\margin{i.e. \uline{si} $V$ a un ``centre sur $U$''\unsure{,} i.e. domine un \unsure{anneau} local de $U$\ldots}
[ \uline{A t-on une réciproque ?} regarder ! ]
\margin{Ça devrait être le cas pour $U$ ``élémentaire'' à facteurs \unsure{pas|par} de \uline{\unsure{types|type} $(0,0)$ ni $(0,1)$\ldots} \unsure{lié à ceci|lié au point} enfin à l'existence d'un $K$}

Quand on dispose d'un modele \uline{projectif} dont $U$ est un ouvert, on peut dire (par définition valuative de la propreté) que $V$ a nécessairement un centre sur $X$, soit $a$, \deleted{et le centre} \inserted{(\uline{NB} $a$ pas} néc. de codim.1\ldots) et la condition envisagée tantot (assurant que $I_{x'}\unsure{\mathrel{\text{---}}|\longrightarrow}\mathcal E_U^{\overline L}$ est nul) est que $a\in U$.
\margin{119}

\scanpage{122}
\unsure{correspondent|correspondant} exactement (via $x\longmapsto\underline{O}_{\mathfrak X,x}$) aux jauges dans $L$ qui sont ent. de type fini sur $K$, de corps des fractions $L$ --- i.e aux $V$ précédents.

Sauf erreur, ces deux types de p$^t$ de $\mathfrak X$ (qui épuisent \deleted{tout} si $n=1$, \unsure{)} \inserted{qui est le seul cas, avec $n=0$, ou $\mathfrak X$ soit un \unsure{schém.|schéma} !)} sont les seuls pour lesquels $\underline{O}_{\mathfrak X,x}$ soit noethérien. Cependant, on a

\margin{$\big|$}
\uline{Proposition} Pour $\forall x\in\mathfrak X$, $\underline{O}_{\mathfrak X,x}$ est un anneau de valuation dans $L$, de corps des fractions $L$, et contenant $K$. L'application $x\longmapsto\underline{O}_{\mathfrak X,x}$ est une bijection de $\mathfrak X$ sur l'ens de ces anneaux dans $L$.
\margin{On pourrait partir d'un schéma excellent $X$, et construire un $\mathfrak X$ relatif au dessus\ldots}

\uline{Dém} a) Soit $f\in L\setminus\underline{O}_{\mathfrak X,x}$, il faut prouver $1/f\in\underline{O}_{\mathfrak X,x}$. Or $f=g/h$, \inserted{$g,h$}$\unsure{\in}\underline{O}_{X,a}$, où $X$ modèle \inserted{\unsure{régulier|regulier}} (propre), $a$ un p$^t$, $f,g$ \deleted{\unsure{premiers} entre eux} --- utilisant la résolution, \deleted{\unsure{on} \uline{sépare} \unsure{doit|deux} les \unsure{diviseurs}} \inserted{OPS $\exists$ système régulier} de paramètres, avec $f,g$ des monômes
\margin{121}
